% GENERATED FILE. Edit canonical lessons, then run npm run generate:books.
% canonical-source-hash: fnv1a64:d6104f42f3104b60
% canonical-lessons: MW-C191-suchna, MW-C191-niyam, MW-C191-mana, MW-C191-savdhan, MW-C191-manno

\chapter{A Notice on the Door}
\label{ch:mw-a-notice-on-the-door}
\hlchaptermodality{191}

\begin{chapteropening}
\textbf{By the end of this chapter:} \emph{I can understand a shop's notice and its rules.}

Complete the last lesson of chapter 191: I can understand a shop's notice and its rules.
\end{chapteropening}

\section[\texorpdfstring{\mw{सूचना} (sūchnā)}{सूचना (sūchnā)}]{\mw{सूचना} (sūchnā) — a notice}

\label{lesson:MW-C191-suchna}



\begin{quote}
Before the new one: say the Marwadi for brave, then the Marwadi for crazy.
\end{quote}

\subsection*{You'll want to know: \mw{सूचना}}
\textbf{\mw{सूचना}} — \emph{sūchnā} — \textquotedblleft{}a notice\textquotedblright{}.

\begin{quote}
\mw{दुकान} \mw{माथै} \mw{नुवी} \mw{सूचना} \mw{है।}
\end{quote}

\begin{quote}
\emph{dukān māthai nuvī sūchnā hai.}
\end{quote}

\begin{quote}
There is a new notice on the shop.
\end{quote}

A feminine noun.

\begin{cousinweb}
\textbf{\mw{सूचना}} (\emph{sūchnā}) is a tatsama, borrowed from Sanskrit \emph{sūcanā}, \textquotedblleft{}indication, communication\textquotedblright{} (Monier-Williams).
\end{cousinweb}

\begin{grammarlens}[title={What you've built}]
One more word for reading notices and messages.
\end{grammarlens}

\subsection*{Your turn}
\begin{itemize}
\raggedright
  \item \emph{Say it:} \emph{sūchnā}
  \item \emph{Say it:} \emph{sūchnā}, once more
  \item \emph{Say it:} \emph{dukān māthai nuvī sūchnā hai.} — There is a new notice on the shop.
  \item \emph{Recall:} say the Marwadi for brave, then the Marwadi for crazy, then say \emph{sūchnā} again
\end{itemize}

\begin{tcolorbox}[breakable,colback=teal!4,colframe=teal!35!black,title={Before you move on}]
What does \mw{सूचना} mean? (\textquotedblleft{}A notice\textquotedblright{}.) Say it once more.
\end{tcolorbox}
\section[\texorpdfstring{\mw{नियम} (niyam)}{नियम (niyam)}]{\mw{नियम} (niyam) — a rule}

\label{lesson:MW-C191-niyam}



\begin{quote}
Before the new one: say the Marwadi for crazy, then the Marwadi for a notice.
\end{quote}

\subsection*{You'll want to know: \mw{नियम}}
\textbf{\mw{नियम}} — \emph{niyam} — \textquotedblleft{}a rule\textquotedblright{}.

\begin{quote}
\mw{ये} \mw{नियम} \mw{घणो} \mw{जरूरी} \mw{है।}
\end{quote}

\begin{quote}
\emph{ye niyam ghaṇo jarūrī hai.}
\end{quote}

\begin{quote}
This rule matters a great deal.
\end{quote}

A masculine noun.

\begin{cousinweb}
\textbf{\mw{नियम}} (\emph{niyam}) is a tatsama, borrowed from Sanskrit \emph{niyama}, \textquotedblleft{}restraining, holding back\textquotedblright{} (Monier-Williams): a rule holds back what it does not allow.
\end{cousinweb}

\begin{grammarlens}[title={What you've built}]
One more word for reading notices and messages.
\end{grammarlens}

\subsection*{Your turn}
\begin{itemize}
\raggedright
  \item \emph{Say it:} \emph{niyam}
  \item \emph{Say it:} \emph{niyam}, once more
  \item \emph{Say it:} \emph{ye niyam ghaṇo jarūrī hai.} — This rule matters a great deal.
  \item \emph{Recall:} say the Marwadi for crazy, then the Marwadi for a notice, then say \emph{niyam} again
\end{itemize}

\begin{tcolorbox}[breakable,colback=teal!4,colframe=teal!35!black,title={Before you move on}]
What does \mw{नियम} mean? (\textquotedblleft{}A rule\textquotedblright{}.) Say it once more.
\end{tcolorbox}
\section[\texorpdfstring{\mw{मना} (manā)}{मना (manā)}]{\mw{मना} (manā) — forbidden}

\label{lesson:MW-C191-mana}



\begin{quote}
Before the new one: say the Marwadi for a notice, then the Marwadi for a rule.
\end{quote}

\subsection*{You'll want to know: \mw{मना}}
\textbf{\mw{मना}} — \emph{manā} — \textquotedblleft{}forbidden\textquotedblright{}.

\begin{quote}
\mw{अठै} \mw{खाणो} \mw{मना} \mw{है।}
\end{quote}

\begin{quote}
\emph{aṭhai khāṇo manā hai.}
\end{quote}

\begin{quote}
No eating here.
\end{quote}

\begin{cousinweb}
\textbf{\mw{मना}} (\emph{manā}) is usually traced, through Persian, to Arabic \emph{manʿ}, \textquotedblleft{}prevention, holding back\textquotedblright{}.
\end{cousinweb}

\begin{grammarlens}[title={What you've built}]
One more word for reading notices and messages.
\end{grammarlens}

\subsection*{Your turn}
\begin{itemize}
\raggedright
  \item \emph{Say it:} \emph{manā}
  \item \emph{Say it:} \emph{manā}, once more
  \item \emph{Say it:} \emph{aṭhai khāṇo manā hai.} — No eating here.
  \item \emph{Recall:} say the Marwadi for a notice, then the Marwadi for a rule, then say \emph{manā} again
\end{itemize}

\begin{tcolorbox}[breakable,colback=teal!4,colframe=teal!35!black,title={Before you move on}]
What does \mw{मना} mean? (\textquotedblleft{}Forbidden\textquotedblright{}.) Say it once more.
\end{tcolorbox}
\section[\texorpdfstring{\mw{सावधान} (sāvdhān)}{सावधान (sāvdhān)}]{\mw{सावधान} (sāvdhān) — careful}

\label{lesson:MW-C191-savdhan}



\begin{quote}
Before the new one: say the Marwadi for a rule, then the Marwadi for forbidden.
\end{quote}

\subsection*{You'll want to know: \mw{सावधान}}
\textbf{\mw{सावधान}} — \emph{sāvdhān} — \textquotedblleft{}careful\textquotedblright{}.

\begin{quote}
\mw{सड़क} \mw{माथै} \mw{सावधान} \mw{रैवो।}
\end{quote}

\begin{quote}
\emph{saṛak māthai sāvdhān raivo.}
\end{quote}

\begin{quote}
Take care on the road.
\end{quote}

\begin{cousinweb}
\textbf{\mw{सावधान}} (\emph{sāvdhān}) is a tatsama, borrowed from Sanskrit \emph{sāvadhāna}, \textquotedblleft{}attentive, heedful\textquotedblright{} (Monier-Williams).
\end{cousinweb}

\begin{grammarlens}[title={What you've built}]
One more word for reading notices and messages.
\end{grammarlens}

\subsection*{Your turn}
\begin{itemize}
\raggedright
  \item \emph{Say it:} \emph{sāvdhān}
  \item \emph{Say it:} \emph{sāvdhān}, once more
  \item \emph{Say it:} \emph{saṛak māthai sāvdhān raivo.} — Take care on the road.
  \item \emph{Recall:} say the Marwadi for a rule, then the Marwadi for forbidden, then say \emph{sāvdhān} again
\end{itemize}

\begin{tcolorbox}[breakable,colback=teal!4,colframe=teal!35!black,title={Before you move on}]
What does \mw{सावधान} mean? (\textquotedblleft{}Careful\textquotedblright{}.) Say it once more.
\end{tcolorbox}
\section[\texorpdfstring{\mw{मानणो} (mānṇo)}{मानणो (mānṇo)}]{\mw{मानणो} (mānṇo) — to obey; to follow (a rule)}

\label{lesson:MW-C191-manno}



\begin{quote}
Before the new one: say the Marwadi for forbidden, then the Marwadi for careful.
\end{quote}

\subsection*{You'll want to know: \mw{मानणो}}
\textbf{\mw{मानणो}} — \emph{mānṇo} — \textquotedblleft{}to obey; to follow (a rule)\textquotedblright{}.

\begin{quote}
\mw{हूं} \mw{हरेक} \mw{नियम} \mw{मानूं} \mw{हूं।}
\end{quote}

\begin{quote}
\emph{hū̃ harek niyam mānū̃ hū̃.}
\end{quote}

\begin{quote}
I follow every rule.
\end{quote}

Past (I …): \mw{म्हैं} \mw{नियम} \mw{मान्यो}, \emph{mhaĩ niyam mānyo}.

\begin{cousinweb}
\textbf{\mw{मानणो}} (\emph{mānṇo}) is inherited from Sanskrit \emph{manyate}, \textquotedblleft{}thinks, honours\textquotedblright{}; CDIAL lists the Old Marwari verb \emph{mānaï}, \textquotedblleft{}to respect, obey\textquotedblright{} (CDIAL 9857).
\end{cousinweb}

\begin{grammarlens}[title={What you've built}]
One more word for reading notices and messages.
\end{grammarlens}

\subsection*{Your turn}
\begin{itemize}
\raggedright
  \item \emph{Say it:} \emph{mānṇo}
  \item \emph{Say it:} \emph{mānṇo}, once more
  \item \emph{Say it:} \emph{hū̃ harek niyam mānū̃ hū̃.} — I follow every rule.
  \item \emph{Recall:} say the Marwadi for forbidden, then the Marwadi for careful, then say \emph{mānṇo} again
\end{itemize}

\begin{tcolorbox}[breakable,colback=teal!4,colframe=teal!35!black,title={Before you move on}]
What does \mw{मानणो} mean? (\textquotedblleft{}To obey; to follow (a rule)\textquotedblright{}.) Say it once more.
\end{tcolorbox}
